\documentclass[aspectratio=169,11pt]{beamer}
\usepackage[UTF8,fontset=none]{ctex}
\setCJKmainfont{Microsoft JhengHei}
\setCJKsansfont{Microsoft JhengHei}
\usepackage{amsmath,amssymb,booktabs,array}
\usepackage{pgfplots}
\pgfplotsset{compat=1.18}
\definecolor{navy}{HTML}{183750}
\definecolor{teal}{HTML}{007E87}
\definecolor{orange}{HTML}{B45A20}
\setbeamercolor{normal text}{fg=navy,bg=white}
\setbeamercolor{structure}{fg=teal}
\setbeamercolor{frametitle}{fg=navy}
\setbeamercolor{alerted text}{fg=orange}
\setbeamerfont{frametitle}{size=\Large,series=\bfseries}
\setbeamertemplate{navigation symbols}{}
\setbeamertemplate{itemize items}[circle]
\setbeamersize{text margin left=8mm,text margin right=8mm}
\setbeamertemplate{footline}{\hspace{8mm}\scriptsize Handout 6 自讀筆記\hfill\insertsectionhead\hfill\insertframenumber/\inserttotalframenumber\hspace{8mm}\vspace{3mm}}
\newcommand{\src}[1]{\par\vfill{\scriptsize\color{teal}#1}}
\newcommand{\slide}[3]{\begin{frame}[t]{#1}#3\src{#2}\end{frame}}
\newcommand{\bi}{\begin{itemize}\setlength{\itemsep}{5pt}}
\newcommand{\ei}{\end{itemize}}
\newcommand{\key}[1]{\par\medskip{\color{teal}\bfseries #1}\par\smallskip}
\newcommand{\dd}{\mathrm{d}}
\newcommand{\bookslide}[3]{\slide{#1}{課本：#2。頁碼為書本印刷頁碼}{#3}}
% Original images extracted losslessly from the supplied textbook PDF.
\newcommand{\originalfigure}[6]{
\begin{frame}[t]{課本原圖 #1：#2}
\centering
\includegraphics[width=.94\linewidth,height=5.05cm,keepaspectratio]{figures/#4}\par
\vspace{2mm}\raggedright\small #5\par
\vfill{\scriptsize\color{teal}原圖：Collinson, 第 3 版，Fig. #1，p. #3。#6}\par
\end{frame}}


\newcommand{\handoutfigure}[3]{
\begin{frame}[t]{講義原頁 #1：#2}
\centering
\includegraphics[width=.97\linewidth,height=5.65cm,keepaspectratio]{figures/#3}\par
\vfill{\scriptsize\color{teal}原圖與文字：賴盈誌老師提供之 Handout，PDF 第 #1 頁。}\par
\end{frame}}
\title{慣性感測、姿態與整合導航}
\subtitle{Handout 6\\繁體中文自讀詳解}
\author{依賴盈誌老師課程講義整理}
\date{CA7528 Digital Avionics\quad Fall 2026}
\begin{document}
\begin{frame}\titlepage\centering\small 概念、原圖、推導、算例與自我檢查\end{frame}
\section{閱讀地圖}
\slide{範圍與閱讀方式}{Handout 6，PDF pp. 1--102}{
\bi
\item 導航座標、INS、陀螺儀、比力、四元數、誤差及融合。
\item 原頁碼一律指講義 PDF 頁序。標示「課本」的頁碼則是 Collinson 第三章印刷頁碼。
\item 先讀概念，再對照原圖。自我檢查先遮住參考解答，再獨立計算或口頭說明。
\item 自編算例與補充推導依題內假設使用。歷史機型、服務與表格依原教材背景理解。
\ei
}
\slide{分次閱讀安排}{Handout 6 閱讀地圖}{
\small
\bi
\item \textbf{導航導讀}：導航狀態與導引命令。
\item \textbf{導航方法}：外部導航與航向參考。
\item \textbf{INS 架構}：平台式與捷聯式 INS。
\item \textbf{狀態與感測}：角速度、比力與初始條件。
\item \textbf{慣性感測器}：MEMS 振動式陀螺儀、RLG、FOG 與 Sagnac 效應。
\item \textbf{加速度計}：質量彈簧與比力量測。
\ei
}
\slide{分次閱讀安排（續）}{Handout 6 閱讀地圖}{
\small
\bi
\item \textbf{誤差與等級}：偏差如何變成導航誤差。
\item \textbf{系統層級}：ISA、IMU 與 INS。
\item \textbf{姿態表示}：Euler 角的奇異性、四元數與單位範數。
\item \textbf{捷聯解算}：機體到 NED 的方向餘弦矩陣。
\item \textbf{地球與重力}：地理緯度與地心緯度、比力、重力與旋轉修正。
\item \textbf{對準與顯示}：位置初始化、對準及姿態顯示。
\ei
}
\slide{分次閱讀安排（續）}{Handout 6 閱讀地圖}{
\small
\bi
\item \textbf{推算導航}：地速分量與經緯度積分。
\item \textbf{整合導航}：外部輔助與 Kalman 濾波概念。
\item \textbf{互補濾波}：頻率分工與時間常數。
\ei
}
\section{導航導讀}
\slide{導航狀態與導引命令}{Handout 6，PDF pp. 1--9；概念整理與補充推導}{
\small
\bi
\item 導航估計位置、速度及姿態，導引根據任務產生路徑或運動指令，控制執行這些指令。
\item 導航解答必須搭配參考座標與時間，單獨一組三維數字沒有完整物理意義。
\item 無線電導航、推算導航與固定參考定位，利用不同資訊建立目前狀態。
\ei
\key{閱讀重點}先定義座標、狀態及初始條件，再選擇感測與解算方法。
}
\handoutfigure{8}{導航狀態與導引命令}{handout06-p008.png}
\handoutfigure{9}{導航狀態與導引命令}{handout06-p009.png}
\slide{自我檢查：導航狀態與導引命令}{依 Handout 範圍自編問題與參考解答}{
\small
\key{先自行作答}
位置、速度與姿態有何不同？為什麼 INS 需要全部三者？
\key{參考解答}
位置描述所在處，速度描述位置隨時間的改變，姿態描述機體方向。INS 需用姿態轉換感測器量測，再積分得到速度與位置。
}
\section{導航方法}
\slide{外部導航與航向參考}{Handout 6，PDF pp. 10--21；概念整理與補充推導}{
\small
\bi
\item 講義介紹地面無線電、Doppler、磁航向等方式，強調不同資訊來源的限制。
\item VOR 提供相對導航台的方位資訊，DME 提供斜距，與單純水平距離不同。
\item RMI 把航向與無線電方位資訊結合顯示，不應簡化成僅靠地磁的指南針。
\item 真北、磁北、航向與地面航跡需分別定義，磁偏角與風都可能影響對照。
\ei
\key{閱讀重點}獨立慣性導航不依賴持續的外部無線電訊號，仍需要初始化與誤差校正。
}
\handoutfigure{11}{外部導航與航向參考}{handout06-p011.png}
\handoutfigure{20}{外部導航與航向參考}{handout06-p020.png}
\handoutfigure{21}{外部導航與航向參考}{handout06-p021.png}
\slide{自我檢查：外部導航與航向參考}{依 Handout 範圍自編問題與參考解答}{
\small
\key{先自行作答}
飛機位於 DME 台正上方，高度差為 3 km，DME 距離是否為零？
\key{參考解答}
不是。理想情況下斜距約為 3 km，即使水平距離為零仍有垂直距離。斜距與水平距離在接近台站時特別需要區分。
}
\section{INS 架構}
\slide{平台式與捷聯式 INS}{Handout 6，PDF pp. 22--27；概念整理與補充推導}{
\small
\bi
\item 平台式以機械穩定平台維持感測器參考方向，捷聯式把感測器直接固定於機體。
\item 捷聯式由角速度與姿態運算建立等效穩定平台，將機體量測轉到導航座標。
\item 捷聯式減少部分機械機構，但增加對計算、校正與數值積分的依賴。
\ei
\key{閱讀重點}兩種架構都需要正確處理姿態、重力與導航參考系的變化。
}
\handoutfigure{24}{平台式與捷聯式 INS}{handout06-p024.png}
\handoutfigure{26}{平台式與捷聯式 INS}{handout06-p026.png}
\handoutfigure{27}{平台式與捷聯式 INS}{handout06-p027.png}
\slide{自我檢查：平台式與捷聯式 INS}{依 Handout 範圍自編問題與參考解答}{
\small
\key{先自行作答}
捷聯式感測器隨機體旋轉，如何知道哪個量測是向北方向？
\key{參考解答}
先用陀螺儀更新姿態，再由姿態建立機體到導航座標的轉換矩陣。量測經矩陣轉換後才得到向北等方向的分量。
}
\section{狀態與感測}
\slide{角速度、比力與初始條件}{Handout 6，PDF pp. 28--34；概念整理與補充推導}{
\small
\bi
\item 陀螺儀量測角運動，加速度計量測比力，兩者都不能單靠單次讀值直接給出完整位置。
\item 姿態需要初始參考與角運動積分，速度與位置也需要初始值。
\item 「完全獨立」描述外部訊號依賴，並不代表長時間無誤差或不需對準。
\ei
\key{閱讀重點}INS 的狀態是在模型與初始條件下，對量測連續積分得到的估計。
}
\handoutfigure{30}{角速度、比力與初始條件}{handout06-p030.png}
\slide{自我檢查：角速度、比力與初始條件}{依 Handout 範圍自編問題與參考解答}{
\small
\key{先自行作答}
若初始速度多估了 $1\,\mathrm{m/s}$，但加速度量測完全正確，60 s 後會有多少位置誤差？
\key{參考解答}
在一維、初始位置正確且速度偏差固定的模型下，$\Delta x=\Delta v_0t=60\,\mathrm m$。正確的加速度無法自動消除未知初始速度偏差。
}
\section{慣性感測器}
\slide{MEMS 振動式陀螺儀}{Handout 6，PDF pp. 35--39；概念整理與補充推導}{
\small
\bi
\item 振動式陀螺儀維持結構振動，旋轉造成耦合方向的感測響應，再轉為角速度訊號。
\item 選用時要檢查偏差、比例因子、雜訊、溫度及量測範圍，不能只看封裝大小。
\item 角速度積分能更新角度，但偏差也會一併積分成姿態誤差。
\ei
\key{閱讀重點}小體積與低成本是系統選型優點，導航精度仍取決於實際性能與校正。
}
\handoutfigure{37}{MEMS 振動式陀螺儀}{handout06-p037.png}
\handoutfigure{38}{MEMS 振動式陀螺儀}{handout06-p038.png}
\slide{自我檢查：MEMS 振動式陀螺儀}{依 Handout 範圍自編問題與參考解答}{
\small
\key{先自行作答}
自編算例：單軸固定 gyro bias 為 $0.1^\circ/\mathrm s$，靜止 120 s 後角度誤差為多少？
\key{參考解答}
若忽略其他誤差與外部校正，$\Delta\theta=b_gt=12^\circ$。多軸導航還會受到姿態耦合與座標轉換影響。
}
\slide{RLG、FOG 與 Sagnac 效應}{Handout 6，PDF pp. 40--48；概念整理與補充推導}{
\small
\bi
\item 光學陀螺儀讓光沿封閉路徑反向傳播，旋轉造成傳播時間或相位差。
\item RLG 以環形雷射腔的相向光束頻差量測旋轉，FOG 以光纖路徑的相位差感測旋轉。
\item 簡化 Sagnac 關係為 $\Delta t=4A\Omega/c^2$，面積向量與旋轉軸的方向須一致。
\item 講義文字提到角加速度時，應注意此基本效應主要對應角速度，且公式不可漏掉旋轉率。
\ei
\key{閱讀重點}量綱核對：面積乘角速度再除以光速平方，結果才是時間。
}
\handoutfigure{43}{RLG、FOG 與 Sagnac 效應}{handout06-p043.png}
\handoutfigure{45}{RLG、FOG 與 Sagnac 效應}{handout06-p045.png}
\handoutfigure{48}{RLG、FOG 與 Sagnac 效應}{handout06-p048.png}
\slide{自我檢查：RLG、FOG 與 Sagnac 效應}{依 Handout 範圍自編問題與參考解答}{
\small
\key{先自行作答}
自編算例：$A=1\,\mathrm{m^2}$、$\Omega=1\,\mathrm{rad/s}$，用 $c=299{,}792{,}458\,\mathrm{m/s}$ 求時間差。
\key{參考解答}
$\Delta t=4/c^2=4.45\times10^{-17}\,\mathrm s$。極小時間差由光學相位或頻率量測反映，多匝 FOG 使用有效包圍面積。
}
\section{加速度計}
\slide{質量彈簧與比力量測}{Handout 6，PDF pp. 49--57；概念整理與補充推導}{
\small
\bi
\item 感測器以測試質量的位移或維持平衡所需的力，推得沿感測軸的比力。
\item 彈簧模型在準靜態下可用 $kx=ma$ 解釋位移與輸入的關係，動態時還需考慮質量與阻尼。
\item 電容、電感等讀出方式把微小位移變成電訊號，需與實際感測原理及校正對照。
\ei
\key{閱讀重點}靜置在桌上的加速度計會感受到支撐作用，自由落體時理想比力近零。
}
\handoutfigure{50}{質量彈簧與比力量測}{handout06-p050.png}
\handoutfigure{56}{質量彈簧與比力量測}{handout06-p056.png}
\handoutfigure{57}{質量彈簧與比力量測}{handout06-p057.png}
\slide{自我檢查：質量彈簧與比力量測}{依 Handout 範圍自編問題與參考解答}{
\small
\key{先自行作答}
自編算例：測試質量 $m=0.01\,\mathrm{kg}$，彈簧剛度 $k=100\,\mathrm{N/m}$，比力大小為 $2\,\mathrm{m/s^2}$，求準靜態位移。
\key{參考解答}
$x=ma/k=0.0002\,\mathrm m=0.2\,\mathrm{mm}$。此為簡化大小關係，正負號須依感測軸及位移方向定義。
}
\section{誤差與等級}
\slide{偏差如何變成導航誤差}{Handout 6，PDF pp. 58--61；概念整理與補充推導}{
\small
\bi
\item 誤差來源包括陀螺儀與加速度計偏差、比例因子、雜訊及軸向不對準。
\item 速度與位置由積分得到，因此持續的小偏差會隨時間累積。
\item 姿態誤差還可能把重力投影到錯誤方向，形成額外水平加速度誤差。
\item 講義的等級、價格與精度表屬背景範例，選型時需用實際設備規格及任務時間。
\ei
\key{閱讀重點}誤差增長不是只有一條固定直線，取決於誤差型態與模型。
}
\handoutfigure{60}{偏差如何變成導航誤差}{handout06-p060.png}
\handoutfigure{61}{偏差如何變成導航誤差}{handout06-p061.png}
\slide{自我檢查：偏差如何變成導航誤差}{依 Handout 範圍自編問題與參考解答}{
\small
\key{先自行作答}
自編算例：一維常值偏差 $b_a=0.02\,\mathrm{m/s^2}$，初始誤差為零，求 60 s 與 300 s 的誤差。
\key{參考解答}
$\Delta v=b_at$、$\Delta x=\tfrac12b_at^2$。60 s 為 $1.2\,\mathrm{m/s}$ 與 36 m。300 s 為 $6\,\mathrm{m/s}$ 與 900 m。此模型忽略姿態與地球效應。
}
\section{系統層級}
\slide{ISA、IMU 與 INS}{Handout 6，PDF pp. 62；概念整理與補充推導}{
\small
\bi
\item 依講義架構，ISA 提供原始慣性資料，IMU 加入感測誤差補償，INS 再加入導航計算。
\item 市場命名可能不同，應以實際介面及規格確認是否輸出原始量測、增量或導航解。
\item 姿態、位置與速度不是所有 IMU 都直接輸出的資訊。
\ei
\key{閱讀重點}設備名稱只能作初步分類，介面定義才決定如何整合。
}
\handoutfigure{62}{ISA、IMU 與 INS}{handout06-p062.png}
\slide{自我檢查：ISA、IMU 與 INS}{依 Handout 範圍自編問題與參考解答}{
\small
\key{先自行作答}
收到一個標為 IMU 的資料包後，至少要核對哪些事項？
\key{參考解答}
需核對角速度或角增量、比力或速度增量、座標軸與正方向、單位、時間標記、更新率、誤差補償及有效性。不能直接把欄位當成位置。
}
\section{姿態表示}
\slide{Euler 角的奇異性}{Handout 6，PDF pp. 63--64；概念整理與補充推導}{
\small
\bi
\item Euler 角用三個有順序的旋轉描述姿態，角速度轉換矩陣依姿態而變。
\item 標準 3--2--1 表示在 $\theta=\pm90^\circ$ 出現 $\tan\theta$ 與 $\sec\theta$ 奇異。
\item 這是參數化限制，並不表示真實剛體姿態不能到達該方向。
\ei
\key{閱讀重點}姿態參數發散與飛機本身失去物理穩定是不同問題。
}
\handoutfigure{64}{Euler 角的奇異性}{handout06-p064.png}
\slide{自我檢查：Euler 角的奇異性}{依 Handout 範圍自編問題與參考解答}{
\small
\key{先自行作答}
飛機做接近垂直的動作，Euler 計算病態，是否代表陀螺儀一定故障？
\key{參考解答}
不代表。量測可能正常，而姿態參數的轉換在奇異點附近敏感。可改用四元數或方向餘弦矩陣運算，再在合適範圍轉成顯示角度。
}
\slide{四元數與單位範數}{Handout 6，PDF pp. 65--69；概念整理與補充推導}{
\small
\bi
\item 純量在前時，$\boldsymbol e=[\cos(\mu/2),\ \hat{\mathbf a}\sin(\mu/2)]^T$，$\hat{\mathbf a}$ 是單位轉軸。
\item 旋轉四元數須滿足 $\boldsymbol e^T\boldsymbol e=1$，數值積分後可正規化。
\item 單位範數保證幾何約束，不能證明姿態估計沒有偏差，也不能證明飛機穩定。
\ei
\key{閱讀重點}四元數的乘法順序與主動／被動旋轉慣例必須一致。
}
\handoutfigure{65}{四元數與單位範數}{handout06-p065.png}
\handoutfigure{67}{四元數與單位範數}{handout06-p067.png}
\slide{自我檢查：四元數與單位範數}{依 Handout 範圍自編問題與參考解答}{
\small
\key{先自行作答}
自編算例：$\tilde{\boldsymbol e}=(0.8,0.1,0.2,0.5)$，求正規化四元數。
\key{參考解答}
範數為 $\sqrt{0.94}=0.969536$，除以此值後約為 $(0.82514,0.10314,0.20628,0.51571)$。正規化修正長度，無法修正錯誤的旋轉方向。
}
\section{捷聯解算}
\slide{機體到 NED 的方向餘弦矩陣}{Handout 6，PDF pp. 70--73；概念整理與補充推導}{
\small
\bi
\item 明確定義 $\mathbf v^n=C_b^n\mathbf v^b$，上標表示座標系，下標表示來源座標。
\item 理想旋轉矩陣滿足 $(C_b^n)^TC_b^n=I$ 與 $\det C_b^n=+1$。
\item 坐標轉換改變向量的分量表示，不改變向量的幾何長度。
\ei
\key{閱讀重點}每個矩陣先寫清楚轉換方向，再代入公式。
}
\handoutfigure{71}{機體到 NED 的方向餘弦矩陣}{handout06-p071.png}
\handoutfigure{72}{機體到 NED 的方向餘弦矩陣}{handout06-p072.png}
\slide{自我檢查：機體到 NED 的方向餘弦矩陣}{依 Handout 範圍自編問題與參考解答}{
\small
\key{先自行作答}
若 $C_b^n=\begin{bmatrix}0&-1&0\\1&0&0\\0&0&1\end{bmatrix}$，將 $(1,0,0)^T$ 從機體轉到 NED。
\key{參考解答}
結果為 $(0,1,0)^T$。機體 x 軸此時對應導航東向。矩陣轉置才是反向轉換，不能任意互換。
}
\section{地球與重力}
\slide{地理緯度與地心緯度}{Handout 6，PDF pp. 74--75；概念整理與補充推導}{
\small
\bi
\item 地球常以略扁的旋轉橢球近似，地理緯度依橢球法線定義，地心緯度依地心連線定義。
\item 講義 p. 74 的 Earth is flat 應避免直譯為地球平坦，需與橢球示意一併理解。
\item 局部 NED 方向隨位置變化，長距離導航不能一直使用同一組局部方向。
\ei
\key{閱讀重點}短距離的局部平面近似與地球整體形狀模型可同時成立於不同尺度。
}
\handoutfigure{75}{地理緯度與地心緯度}{handout06-p075.png}
\slide{自我檢查：地理緯度與地心緯度}{依 Handout 範圍自編問題與參考解答}{
\small
\key{先自行作答}
為什麼高緯度地區同樣的經度差對應較短的東西向距離？
\key{參考解答}
球形近似下，緯線圈半徑為 $R\cos\varphi$，東向小距離約 $R\cos\varphi\,\Delta\lambda$。越靠近極區，$\cos\varphi$ 越小。
}
\slide{比力、重力與旋轉修正}{Handout 6，PDF pp. 76--80；概念整理與補充推導}{
\small
\bi
\item 在局部平面且忽略地球自轉的簡化下，$\dot{\mathbf v}^n=C_b^n\mathbf f^b+\mathbf g^n$。
\item 更完整的導航方程還要處理地球自轉、導航座標轉動與 Coriolis 項。
\item 講義有不同字母與軸定義，抄公式前要先對齊符號及重力正方向。
\ei
\key{閱讀重點}直接積分未轉軸、未加回重力的加速度計讀值，通常得不到正確導航速度。
}
\handoutfigure{76}{比力、重力與旋轉修正}{handout06-p076.png}
\handoutfigure{79}{比力、重力與旋轉修正}{handout06-p079.png}
\slide{自我檢查：比力、重力與旋轉修正}{依 Handout 範圍自編問題與參考解答}{
\small
\key{先自行作答}
NED 下靜置且機體與導航軸一致，$\mathbf f^b=(0,0,-9.81)^T$、$\mathbf g^n=(0,0,9.81)^T$，求運動加速度。
\key{參考解答}
兩者相加為零，符合靜止條件。這也說明加速度計輸出與位置二階導數之間，還需要重力與參考系的處理。
}
\section{對準與顯示}
\slide{位置初始化、對準及姿態顯示}{Handout 6，PDF pp. 81--88；概念整理與補充推導}{
\small
\bi
\item 初始化包含位置與速度等狀態，對準建立初始姿態及感測軸相對導航系的方向。
\item 陀螺儀與加速度計誤差、飛機運動及外部參考品質都可能影響對準。
\item 講義的對準時間屬設備與方法背景，不能當成所有 INS 的固定規格。
\ei
\key{閱讀重點}顯示器呈現的是估計結果，對準與有效性狀態應納入判讀。
}
\handoutfigure{83}{位置初始化、對準及姿態顯示}{handout06-p083.png}
\handoutfigure{87}{位置初始化、對準及姿態顯示}{handout06-p087.png}
\slide{自我檢查：位置初始化、對準及姿態顯示}{依 Handout 範圍自編問題與參考解答}{
\small
\key{先自行作答}
姿態畫面平穩且數值不跳動，能否證明位置導航也準確？
\key{參考解答}
不能。平穩可能只是雜訊低，位置仍會受到初始位置、速度及感測偏差累積影響。需使用獨立參考與誤差評估。
}
\section{推算導航}
\slide{地速分量與經緯度積分}{Handout 6，PDF pp. 89--93；概念整理與補充推導}{
\small
\bi
\item 短距離球形近似為 $\dot\varphi=v_N/R$、$\dot\lambda=v_E/(R\cos\varphi)$，角度採 rad。
\item 速度須為相對地面的北、東分量。空速與機頭航向還需要考慮風才能得到地面運動。
\item 公式在極區需特別處理，長距離也需要更新曲率半徑與緯度。
\ei
\key{閱讀重點}計算前先把角度轉成 rad，更新完成後再按需要轉成度。
}
\handoutfigure{90}{地速分量與經緯度積分}{handout06-p090.png}
\handoutfigure{92}{地速分量與經緯度積分}{handout06-p092.png}
\slide{自我檢查：地速分量與經緯度積分}{依 Handout 範圍自編問題與參考解答}{
\small
\key{先自行作答}
自編例：$\varphi_0=25^\circ$、$\lambda_0=121^\circ$、$v_N=120$、$v_E=80\,\mathrm{m/s}$，$t=60\,\mathrm s$、$R=6.371\times10^6\,\mathrm m$。
\key{參考解答}
用初始緯度近似，$\Delta\varphi=0.06475^\circ$、$\Delta\lambda=0.04763^\circ$。新位置約為 $(25.06475^\circ,121.04763^\circ)$。
}
\section{整合導航}
\slide{外部輔助與 Kalman 濾波概念}{Handout 6，PDF pp. 94--97；概念整理與補充推導}{
\small
\bi
\item INS 短時間反應快且連續，外部位置或速度資訊可限制其長時間漂移。
\item Kalman 濾波使用狀態模型預測，再依量測殘差與不確定度更新估計。
\item 量測融合前仍需同步時間、轉換座標與檢查離群值，權重不是固定偏好某一種感測器。
\ei
\key{閱讀重點}融合成效依賴模型、量測品質與有效的故障偵測。
}
\handoutfigure{95}{外部輔助與 Kalman 濾波概念}{handout06-p095.png}
\handoutfigure{97}{外部輔助與 Kalman 濾波概念}{handout06-p097.png}
\slide{自我檢查：外部輔助與 Kalman 濾波概念}{依 Handout 範圍自編問題與參考解答}{
\small
\key{先自行作答}
GPS 中斷後，整合 INS 是否立刻沒有導航輸出？
\key{參考解答}
通常仍可依 INS 傳播狀態，但缺少外部修正後不確定度會增長。可維持多久取決於誤差、動態與任務要求，不能給所有設備同一時限。
}
\section{互補濾波}
\slide{頻率分工與時間常數}{Handout 6，PDF pp. 98--102；概念整理與補充推導}{
\small
\bi
\item 令 $F_1(s)=Ts/(1+Ts)$、$F_2(s)=1/(1+Ts)$，則 $F_1+F_2=1$。
\item 快速、低雜訊但會漂移的來源走高通，長期準確但較吵的來源走低通。
\item 交越量級為 $\omega_c=1/T$，此時不是兩個實數各等於一半，還有相位差。
\ei
\key{閱讀重點}互補是傳遞函數相加為一，並不代表任何頻率都能消除全部誤差。
}
\handoutfigure{100}{頻率分工與時間常數}{handout06-p100.png}
\handoutfigure{101}{頻率分工與時間常數}{handout06-p101.png}
\handoutfigure{102}{頻率分工與時間常數}{handout06-p102.png}
\slide{自我檢查：頻率分工與時間常數}{依 Handout 範圍自編問題與參考解答}{
\small
\key{先自行作答}
自編算例：$T=5\,\mathrm s$，求低頻、高頻極限及交越角頻率。
\key{參考解答}
$s\to0$ 時 $F_1\to0,F_2\to1$。高頻時 $F_1\to1,F_2\to0$。$\omega_c=0.2\,\mathrm{rad/s}$，兩者幅值在此均為 $1/\sqrt2$。
}
\section{複習與來源}
\slide{整合複述檢查}{Handout 6 全份複習}{
\bi
\item 選一項系統功能，說明輸入、輸出、座標、單位與資料來源。
\item 找出一個公式，逐一解釋每個符號，再說明它使用哪些假設。
\item 選一個故障或誤差，追蹤它如何影響下游功能，以及可以用什麼獨立資訊核對。
\item 重做本份的自我檢查，再對照講義與相關公式。
\ei
}
\slide{來源與後續閱讀}{Handout 6 與隨附教材}{
\small
\bi
\item 主要來源：賴盈誌老師，CA7528 Digital Avionics，Handout 6。原始檔保留於上一層資料夾。
\item 版型與重用內容：既有 Chapter03\_SelfStudy\_LaTeX。Handout 3、4 沿用相關章節及其原圖。
\item 第三章重用內容來源：R. P. G. Collinson，\textit{Introduction to Avionics Systems}，第 3 版。
\item 本筆記的重算、核對與延伸為自讀整理，並非授課教師或出版社正式勘誤。
\item 原圖內來源標示與圖說保留，新增題目與參考解答已另行標明。
\ei
}
\end{document}
